\documentclass[multi=page, border=15pt, varwidth=16cm]{standalone}
\usepackage{ctex}      % 提供中文支持
\usepackage{tabularx}  % 自动调整列宽
\usepackage{booktabs}  % 优雅的三线表支持
\usepackage{xcolor}
\usepackage{colortbl}

% 统一的深灰蓝表头色
\definecolor{headerbg}{RGB}{52, 73, 94}
\definecolor{headertext}{RGB}{255, 255, 255}

% 自定义表头文本样式命令
\newcommand{\thnode}[1]{\textcolor{headertext}{\bfseries #1}}

\begin{document}
\sffamily
\renewcommand{\arraystretch}{1.5}

% ================= 图片：fdisk 交互指令集 =================
\begin{page}
	\centering
	% \textbf{\Large fdisk 磁盘分区交互指令集} \\ \vspace{0.8em}

	\begin{tabularx}{\linewidth}{>{\bfseries\color{blue!70!black}}c X >{\bfseries\color{blue!70!black}}c X >{\bfseries\color{blue!70!black}}c X}
		\toprule
		\rowcolor{headerbg}
		\thnode{指令} & \thnode{说明}      & \thnode{指令} & \thnode{说明}    & \thnode{指令} & \thnode{说明} \\
		\midrule
		% 第一行：最核心的查看与新建
		p           & \textbf{显示当前分区表} & n           & \textbf{添加新分区} & m           & 显示帮助菜单      \\
		% 第二行：保存与退出
		w           & \textbf{保存更改并退出} & q           & 不保存直接退出        & d           & 删除现有分区      \\
		% 第三行：类型与验证
		l           & 列出已知分区类型         & t           & 修改分区系统 ID      & v           & 验证分区表内容     \\
		% 第四行：特殊与高级
		a           & 设置活动分区标志         & u           & 切换显示单位         & x           & 进入高级专家模式    \\
		\bottomrule
		\multicolumn{6}{p{\linewidth}}{\small \color{gray}\textit{* 注：在交互界面按下 \texttt{n} 后，通常需要进一步选择 \texttt{p}(主分区) 或 \texttt{e}(扩展分区)，并指定大小（如 \texttt{+2G}）}}
	\end{tabularx}

	% \vspace{1.5em}

	% \begin{tabularx}{\linewidth}{l X}
	% 	\rowcolor{headerbg!10}
	% 	\textbf{核心操作流：} & \texttt{fdisk /dev/sdb} $\rightarrow$ \texttt{p} (查看) $\rightarrow$ \texttt{n} (新建) $\rightarrow$ \texttt{w} (写入生效) \\
	% \end{tabularx}
\end{page}

\end{document}